\section{正则化、岭回归与 Lasso}\label{ch:regularization}

\subsection{为什么需要限制模型的复杂度}

正则化为模型选择增加了一种归纳偏好：当若干候选模型的训练拟合表现相近时，优先考虑某些结构更简单或参数更稳定的模型。简单的具体含义取决于惩罚函数，例如岭回归偏好幅度较小的系数，Lasso 则可能产生精确为零的系数。这样的偏好可能牺牲部分训练拟合效果，却有机会降低模型对有限样本中偶然波动的敏感程度。

第\ref{ch:generalization}节说明：模型在训练样本上误差很小，并不保证在新样本上预测准确。例如，给定学习时间 $x=1,2,3,4,5$ 与成绩 $y=60,65,70,75,80$，直线 $\hat y=5x+55$ 可以完全拟合训练数据。但若采用次数更高的多项式，也可能达到同样的零训练误差，却在未见数据上作出很不相同的预测。因此，除了拟合训练数据，我们还希望模型不要无谓地复杂。

一种常见的建模策略是\textbf{正则化}（Regularization）：在原有训练损失之外加入对参数复杂程度的惩罚。它不是简单地认定“系数越小，模型越正确”，而是在拟合误差和模型复杂程度之间做出明确的权衡\cite[第11章]{zhou2016}。

假设两个候选模型的训练误差相等，系数向量分别为
\[
\boldsymbol w_A=(1,2)^{\mathsf T},\qquad
\boldsymbol w_B=(5,10)^{\mathsf T}.
\]
若使用系数平方和衡量复杂程度，则 $\|\boldsymbol w_A\|_2^2=1^2+2^2=5$，而 $\|\boldsymbol w_B\|_2^2=5^2+10^2=125$。在这一特定准则下，模型 A 所受惩罚较小。但不同特征的量纲会影响系数数值，所以实际比较之前通常需要适当标准化特征。

\subsection{从普通最小二乘到正则化目标}

加入惩罚项后，优化目标已经改变，因此不能再要求正则化解同时使原残差平方和最小。模型寻找的是拟合与惩罚之间的平衡。正则化强度由建模者选择，而不是作为普通权重随训练目标一起自由最小化；否则算法往往倾向于取消惩罚，从而失去约束模型复杂度的目的。

沿用第\ref{ch:least-squares}节的单变量线性回归 $\hat y_i=wx_i+b$，原来的训练目标是
\[
J_{\mathrm{OLS}}(w,b)=\frac1{2n}\sum_{i=1}^{n}(y_i-wx_i-b)^2.
\]
现在增加一个系数平方惩罚：
\begin{equation}\label{eq:ridge-scalar}
\boxed{J_{\mathrm{Ridge}}(w,b)=
\underbrace{\frac1{2n}\sum_{i=1}^{n}(y_i-wx_i-b)^2}_{\text{数据拟合项}}
+\underbrace{\frac{\lambda}{2}w^2}_{\text{正则化项}},\qquad\lambda\geq0.}
\end{equation}
$\lambda$ 是事先选择的\textbf{超参数}，决定惩罚的强度；$w,b$ 才是通过训练求解的参数。这里沿用通常不惩罚截距 $b$ 的约定。当 $\lambda=0$ 时退回普通最小二乘；在数据项保持不变的比较中，$\lambda$ 越大，大系数所付的代价越高。

例如两个样本的真实值为 $5,8$，预测值为 $4,6$，则数据拟合项为
\[
J_{\mathrm{error}}=\frac1{2\times2}\big[(5-4)^2+(8-6)^2\big]=1.25.
\]
若这组参数满足 $w=4$ 且选择 $\lambda=1$，则总目标为
\[
J_{\mathrm{Ridge}}=1.25+\frac12\times4^2=\boxed{9.25}.
\]
\textbf{注意：}该数值是包含惩罚项的训练目标，不等于 MSE，也不能直接当作测试误差。

\subsection{手工推导岭回归的参数收缩}

考虑已中心化且无截距的两个教学样本 $(x_1,y_1)=(1,3)$、$(x_2,y_2)=(-1,-3)$，模型取 $\hat y=wx$。普通最小二乘的最优参数为 $w_{\mathrm{OLS}}=3$。加入正则化后，
\begin{align}
J(w)
&=\frac14\big[(3-w)^2+(-3+w)^2\big]+\frac{\lambda}{2}w^2\notag\\
&=\frac12(3-w)^2+\frac{\lambda}{2}w^2.\label{eq:ridge-hand-objective}
\end{align}
先求导：
\[
J'(w)=-(3-w)+\lambda w=(1+\lambda)w-3.
\]
令导数为零，得到
\begin{equation}\label{eq:ridge-hand-solution}
\boxed{w^*=\frac{3}{1+\lambda}.}
\end{equation}
由于 $J''(w)=1+\lambda>0$，这一驻点是唯一的全局最小点。代入不同的正则化强度，得到
\begin{center}
\begin{tabular}{rrrrrr}\toprule
$\lambda$&0&0.5&1&2&5\\\midrule
$w^*$&3&2&1.5&1&0.5\\\bottomrule
\end{tabular}
\end{center}
在\emph{这个单变量例子}中，$\lambda$ 增大会使系数逐渐接近零。不能由此断言在一般多特征岭回归中，每一个系数的绝对值都必然随 $\lambda$ 单调减小；当特征相关时，系数可能相互补偿。

\subsection{多特征岭回归及其矩阵形式}

岭回归尤其能够帮助理解多重共线性。如果两个特征高度相关，普通最小二乘可能产生幅度很大且方向相反的系数，虽然总体预测尚可，单独系数却容易因数据的小幅扰动发生变化。岭惩罚不鼓励这种大幅度抵消，从而提高数值稳定性。对于中心化设计矩阵，正的岭惩罚系数使系数矩阵获得正定的对角修正，即使原设计矩阵不满列秩，也能唯一确定受惩罚的权重。

设有 $p$ 个数值特征，参数向量 $\boldsymbol w\in\mathbb R^p$，截距为 $b$。预测模型为 $\hat y_i=\boldsymbol w^{\mathsf T}\boldsymbol x_i+b$。定义
\[
\|\boldsymbol w\|_2^2=w_1^2+w_2^2+\cdots+w_p^2.
\]
于是\textbf{岭回归}（Ridge Regression）的目标函数为
\begin{equation}\label{eq:ridge-multivariate}
\boxed{J_{\mathrm{Ridge}}(\boldsymbol w,b)=
\frac1{2n}\sum_{i=1}^{n}
(y_i-\boldsymbol w^{\mathsf T}\boldsymbol x_i-b)^2+
\frac{\lambda}{2}\|\boldsymbol w\|_2^2.}
\end{equation}
若数据已经中心化且不另设截距，设计矩阵为 $X\in\mathbb R^{n\times p}$，则
\[
J(\boldsymbol w)=\frac1{2n}\|X\boldsymbol w-\boldsymbol y\|_2^2+
\frac{\lambda}{2}\|\boldsymbol w\|_2^2.
\]
对 $\boldsymbol w$ 求梯度：
\[
\nabla J(\boldsymbol w)=\frac1nX^{\mathsf T}(X\boldsymbol w-\boldsymbol y)+\lambda\boldsymbol w.
\]
令梯度为零，整理得
\begin{equation}\label{eq:ridge-normal}
(X^{\mathsf T}X+n\lambda I_p)\hat{\boldsymbol w}=X^{\mathsf T}\boldsymbol y.
\end{equation}
当 $\lambda>0$ 时，$X^{\mathsf T}X+n\lambda I_p$ 正定，因此该线性方程组有唯一解，即使 $X$ 本身不满列秩也成立。这里的 $I_p$ 是 $p$ 阶单位矩阵；若要保留不受惩罚的截距，应先中心化或在惩罚矩阵中单独处理截距位置。

\subsection{Lasso：为什么会出现零系数}

Lasso 与岭回归的关键区别来自惩罚函数在零附近的几何性质。平方惩罚在零点的斜率为零，通常使参数连续缩小；绝对值惩罚在零处有折角，一定范围内的拟合梯度可以被零点的次梯度抵消，使最优参数恰好为零。零系数只表示特征没有进入当前预测函数，不能据此认定它在总体中没有任何作用。相关特征之间的选择可能不稳定，变量单位不同也会影响惩罚的相对强度。

\textbf{Lasso} 将平方惩罚换为绝对值惩罚：
\begin{equation}\label{eq:lasso-objective}
\boxed{J_{\mathrm{Lasso}}(\boldsymbol w,b)=
\frac1{2n}\sum_{i=1}^{n}(y_i-\boldsymbol w^{\mathsf T}\boldsymbol x_i-b)^2+
\lambda\|\boldsymbol w\|_1,\quad
\|\boldsymbol w\|_1=\sum_{j=1}^{p}|w_j|.}
\end{equation}
例如 $\boldsymbol w=(1,-2,3)^{\mathsf T}$，则 $\|\boldsymbol w\|_2^2=14$，而 $\|\boldsymbol w\|_1=6$。L2 与 L1 惩罚所用的量并不相同，不能只比较这两个数值的大小来判定哪个算法更好。

\begin{center}
\begin{tabularx}{\textwidth}{p{28mm}XX}\toprule
比较项目&岭回归（Ridge）&Lasso\\\midrule
惩罚形式&系数平方和（L2）&系数绝对值和（L1）\\
参数特点&系数通常缩小但未必为零&某些系数可恰好等于零\\
典型用途&稳定估计、处理相关特征&稀疏建模、嵌入式特征选择\\
数学性质&光滑凸目标&凸但在系数零处不可微\\\bottomrule
\end{tabularx}
\end{center}

为什么 Lasso 可能出现精确的零？仍以先前的单参数样本为例，其目标为
\[
J_{\mathrm{Lasso}}(w)=\frac12(3-w)^2+\lambda|w|.
\]
当 $w>0$ 时，一阶条件是 $w-3+\lambda=0$，得到 $w=3-\lambda$；当 $w<0$ 时，一阶条件给出 $w=3+\lambda>0$，与 $w<0$ 矛盾。结合 $w=0$ 处的最优性条件，可得
\begin{equation}\label{eq:lasso-soft-example}
\boxed{w_{\mathrm{Lasso}}^*=\max(3-\lambda,0),\qquad\lambda\geq0.}
\end{equation}
因此当 $\lambda\geq3$ 时系数恰好为零。这个单参数推导说明 L1 惩罚具有产生稀疏解的机制；多特征问题的求解需要考虑各参数及特征的相互作用。

\subsection{如何选择正则化强度}

并不是 $\lambda$ 越大越好。惩罚过弱时，模型可能仍然过拟合；惩罚过强时，参数可能被过度压缩，出现欠拟合。实际建模通常在训练集内部使用交叉验证选择 $\lambda$，并在最终独立测试集上评价模型\cite[第5章]{james2023}。如需标准化特征，必须仅用每个训练折估计均值和标准差，以免产生第\ref{ch:generalization}节讨论的数据泄漏。

\subsection{本章小结}

普通最小二乘只最小化数据拟合误差；岭回归额外惩罚系数平方和；Lasso 则惩罚系数绝对值之和。它们体现了统一的统计学习思想：\textbf{优化目标中的“拟合得好”与“模型不要过于复杂”可以同时考虑。} 其中，正则化优化问题求得最优解并不等于保证模型具有最好的泛化能力。
